\textbf{Present Value of an Annuity}

\begin{definition}
    The \emph{present value of an annuity} is the sum of the present values of all payments subject to some fixed rate.
\end{definition}

The present value of an annuity represents the amount that must be invested now at rate \(r\) so that all payments can be made.
Recall from previous sections that the present value of each payment, \(R,\) subject to the periodic rate \(r\) can be calculated as in \autoref{presentValueFixedPayment}. 

\begin{table}[]
    \centering
    \begin{tabular}{c|c}
         Payment number & Present Value \\ \hline
         \(1\)st & \(R\left(1+r\right)^{-1}\) \\
         \(2\)nd & \(R\left(1+r\right)^{-2}\) \\
         \(3\)rd & \(R\left(1+r\right)^{-3}\) \\
         \vdots & \vdots \\
         \(n\)th & \(R\left(1+r\right)^{-n}\) \\
    \end{tabular}
    \caption{Present value of each of the \(n\) payments}
    \label{presentValueFixedPayment}
\end{table}

The present value of the \emph{annuity}, \(PV_a,\) is the sum of the present values of each payment. From \autoref{presentValueFixedPayment} we see that 

\[PV_a = R\left(1+r\right)^{-1} + R\left(1+r\right)^{-2} + \dots + R\left(1+r\right)^{-n}.\]

This is the sum of a geometric sequence and hence the formula that we have for sums of geometric sequences applies.

\begin{definition}
    The formula
    \[PV_a = R\frac{1-\left(1+r\right)^{-n}}{r}\]
    gives the \emph{present value} \(PV_a\) of an ordinary annuity of \(R\) per payment period for \(n\) periods at the periodic interest rate of \(r.\)
\end{definition}

It will be helpful for us to represent this as a function of \(n\) and \(r.\) Define \(a\left(n,r\right) = \frac{1-\left(1+r\right)^{-n}}{r}.\)